%%
%% This is file `sample-sigconf-authordraft.tex',
%% generated with the docstrip utility.
%%
%% The original source files were:
%%
%% samples.dtx  (with options: `all,proceedings,bibtex,authordraft')
%% 
%% IMPORTANT NOTICE:
%% 
%% For the copyright see the source file.
%% 
%% Any modified versions of this file must be renamed
%% with new filenames distinct from sample-sigconf-authordraft.tex.
%% 
%% For distribution of the original source see the terms
%% for copying and modification in the file samples.dtx.
%% 
%% This generated file may be distributed as long as the
%% original source files, as listed above, are part of the
%% same distribution. (The sources need not necessarily be
%% in the same archive or directory.)
%%
%%
%% Commands for TeXCount
%TC:macro \cite [option:text,text]
%TC:macro \citep [option:text,text]
%TC:macro \citet [option:text,text]
%TC:envir table 0 1
%TC:envir table* 0 1
%TC:envir tabular [ignore] word
%TC:envir displaymath 0 word
%TC:envir math 0 word
%TC:envir comment 0 0
%%
%% The first command in your LaTeX source must be the \documentclass
%% command.
%%
%% For submission and review of your manuscript please change the
%% command to \documentclass[manuscript, screen, review]{acmart}.
%%
%% When submitting camera ready or to TAPS, please change the command
%% to \documentclass[sigconf]{acmart} or whichever template is required
%% for your publication.
%%
%%
\documentclass[sigconf,authorversion,nonacm]{acmart}
%%
%% \BibTeX command to typeset BibTeX logo in the docs
\AtBeginDocument{%
  \providecommand\BibTeX{{%
    Bib\TeX}}}
\usepackage{arydshln}
\usepackage{nicematrix}

%% Rights management information.  This information is sent to you
%% when you complete the rights form.  These commands have SAMPLE
%% values in them; it is your responsibility as an author to replace
%% the commands and values with those provided to you when you
%% complete the rights form.
% \setcopyright{acmlicensed}
% \copyrightyear{2018}
% \acmYear{2018}
% \acmDOI{XXXXXXX.XXXXXXX}
% %% These commands are for a PROCEEDINGS abstract or paper.
% \acmConference[Conference acronym 'XX]{Make sure to enter the correct
%   conference title from your rights confirmation email}{June 03--05,
%   2018}{Woodstock, NY}
% %%
% %%  Uncomment \acmBooktitle if the title of the proceedings is different
% %%  from ``Proceedings of ...''!
% %%
% %%\acmBooktitle{Woodstock '18: ACM Symposium on Neural Gaze Detection,
% %%  June 03--05, 2018, Woodstock, NY}
% \acmISBN{978-1-4503-XXXX-X/2018/06}


%%
%% Submission ID.
%% Use this when submitting an article to a sponsored event. You'll
%% receive a unique submission ID from the organizers
%% of the event, and this ID should be used as the parameter to this command.
%%\acmSubmissionID{123-A56-BU3}

%%
%% For managing citations, it is recommended to use bibliography
%% files in BibTeX format.
%%
%% You can then either use BibTeX with the ACM-Reference-Format style,
%% or BibLaTeX with the acmnumeric or acmauthoryear sytles, that include
%% support for advanced citation of software artefact from the
%% biblatex-software package, also separately available on CTAN.
%%
%% Look at the sample-*-biblatex.tex files for templates showcasing
%% the biblatex styles.
%%

%%
%% The majority of ACM publications use numbered citations and
%% references.  The command \citestyle{authoryear} switches to the
%% "author year" style.
%%
%% If you are preparing content for an event
%% sponsored by ACM SIGGRAPH, you must use the "author year" style of
%% citations and references.
%% Uncommenting
%% the next command will enable that style.
%%\citestyle{acmauthoryear}
\newcommand{\bai}[1]{{
 		{\color{red}\noindent\textbf{BC:}
 			\emph{#1}
 		}\normalcolor}}

%%
%% end of the preamble, start of the body of the document source.
\begin{document}

%%
%% The "title" command has an optional parameter,
%% allowing the author to define a "short title" to be used in page headers.
\title{Ensemble Multi-Objective Optimization on MOOT Datasets \vspace{0.5em} \\ \large CSC591 Team 8 \vspace{1.5em}}

%%
%% The "author" command and its associated commands are used to define
%% the authors and their affiliations.
%% Of note is the shared affiliation of the first two authors, and the
%% "authornote" and "authornotemark" commands
%% used to denote shared contribution to the research.

\author{Bhavana Chappidi}
\affiliation{
  \institution{}
  \city{}
  \country{}
 }
\email{bchappi@ncsu.edu}

\author{Nishad Tardalkar}
\affiliation{
  \institution{}
  \city{}
  \country{}
}
\email{ntardal@ncsu.edu}



%%
%% By default, the full list of authors will be used in the page
%% headers. Often, this list is too long, and will overlap
%% other information printed in the page headers. This command allows
%% the author to define a more concise list
%% of authors' names for this purpose.
\renewcommand{\shortauthors}{P08 CSC 591 SWEAI Project}


\maketitle


\settopmatter{printacmref=false}
\renewcommand\footnotetextcopyrightpermission[1]{}



\section{Introduction}

Many real-world problems in software engineering involve balancing multiple 
conflicting objectives. For example, a system configuration may need to 
maximize performance while minimizing energy consumption, memory usage, 
or operational cost. Similarly, machine learning systems often require 
trade-offs between model accuracy, training time, and computational 
resources. These types of problems cannot be effectively solved using 
single-objective optimization techniques because improving one objective 
often leads to worsening another.

Multi-objective optimization (MOO) provides a framework for solving such 
problems. Instead of producing a single optimal solution, MOO algorithms 
generate a set of trade-off solutions known as the Pareto frontier. Each 
solution on the Pareto frontier represents a different compromise among 
objectives, and none of these solutions is strictly better than another 
across all objectives.

Over the past two decades, many multi-objective optimization algorithms 
have been developed. Examples include NSGA-II~\cite{deb2002nsga2}, MOEA/D~\cite{zhang2007moead}, 
and other evolutionary algorithms that explore large search spaces and attempt 
to approximate the Pareto frontier. Frameworks such as \textit{pymoo}~\cite{blank2020pymoo} 
provide implementations of these algorithms and allow researchers to evaluate them 
on a variety of optimization problems.

Despite the large number of available algorithms, most prior work 
evaluates them individually. Researchers typically select a single 
algorithm, run it on a dataset, and analyze the results. However, 
different algorithms often explore the search space in different ways. 
As a result, each algorithm may discover different regions of the Pareto 
frontier.

In machine learning and predictive modeling, combining multiple models 
into an ensemble has often been shown to improve performance. Ensemble methods leverage the strengths of multiple models to produce more robust 
and accurate results than any single model alone~\cite{dietterich2000ensemble}. 
A similar idea may be applicable to multi-objective optimization.

This project investigates whether combining the final Pareto frontiers 
produced by multiple optimization algorithms can lead to better 
approximations of the true Pareto frontier. Specifically, we examine 
whether pooling the solutions found by different algorithms and removing 
dominated solutions can create a stronger overall frontier than the 
results produced by any individual algorithm.

\section{Research Questions}

The primary goal of this study is to explore whether ensemble techniques 
can improve multi-objective optimization results in the context of 
software engineering datasets.

The main research question is:

\textbf{RQ:} Can an ensemble of multi-objective optimization algorithms, 
formed by pooling their final Pareto frontiers, outperform individual 
algorithms on software engineering datasets?

To better understand this question, the study also investigates several 
supporting questions:

\begin{itemize}
\item Do different optimization algorithms discover different regions 
of the Pareto frontier?
\item Does combining Pareto fronts from multiple algorithms improve 
performance metrics such as hypervolume and spread?
\item Can solutions obtained from an ensemble Pareto frontier improve 
optimization performance when used as initialization for another algorithm?
\end{itemize}

\section{Contributions}

This study makes several contributions to the study of search-based 
software engineering and multi-objective optimization:

\begin{itemize}
\item An empirical comparison of several multi-objective optimization 
algorithms on software engineering datasets.

\item An investigation into whether combining the final Pareto frontiers 
of multiple algorithms can produce higher quality solutions than those 
generated by individual algorithms.

\item An experimental evaluation of ensemble-based initialization, 
where solutions obtained from combined Pareto fronts are used to seed 
future optimization runs.

\item A reproducible experimental framework implemented using the 
pymoo optimization library and datasets from the MOOT repository~\cite{menzies2025moot}.
\end{itemize}

\section{Background}

\subsection{Literature Overview}

Multi-objective optimization has been widely studied across a variety of domains, 
including engineering design, energy systems, and machine learning. Prior work 
has focused on developing optimization algorithms, benchmarking their performance, 
and analyzing evaluation metrics. However, most studies evaluate algorithms 
individually rather than considering how their outputs might be combined.

\subsection{Multi-Objective Optimization}

Multi-objective optimization refers to problems that contain multiple 
objectives that must be optimized simultaneously. In many real-world 
scenarios, these objectives conflict with one another. For example, 
improving system performance may require additional hardware resources, 
which increases cost.

Instead of searching for a single optimal solution, multi-objective 
optimization algorithms aim to identify a set of solutions that represent 
different trade-offs between objectives. These solutions form what is 
known as the Pareto frontier.

\subsection{Pareto Dominance}

A solution is considered \textit{Pareto dominated} if another solution 
exists that performs at least as well in all objectives and strictly 
better in at least one objective. Solutions that are not dominated by 
any other solutions are called \textit{non-dominated}. The collection of 
these non-dominated solutions forms the Pareto frontier.

Pareto frontiers provide decision-makers with a range of choices. 
Instead of relying on a single optimal configuration, practitioners 
can examine multiple trade-off solutions and select the one that best 
matches their priorities.

\subsection{Evolutionary Multi-Objective Algorithms}

Many multi-objective optimization algorithms are based on evolutionary 
search techniques. These algorithms maintain a population of candidate 
solutions and iteratively improve them using operations such as mutation, 
selection, and crossover.

Algorithms such as NSGA-II~\cite{deb2002nsga2} and MOEA/D~\cite{zhang2007moead} 
are widely used because they maintain diverse populations and encourage 
exploration of different regions of the search space. These algorithms attempt to approximate the 
true Pareto frontier by evolving solutions over multiple generations.


\subsection{Evaluation Metrics}

Several metrics have been proposed to evaluate the quality of Pareto 
frontiers~\cite{riquelme2015metrics,audet2021performance}. In this study, 
we use four metrics to evaluate both convergence and diversity of solutions.

\textbf{Hypervolume (HV)} measures the volume of objective space dominated 
by the solutions in the Pareto frontier~\cite{shang2020hypervolume}. A larger 
hypervolume indicates better coverage of the objective space and closer 
approximation to the optimal frontier.

\textbf{Inverted Generational Distance (IGD)} measures the average distance 
from a set of reference Pareto-optimal solutions to the obtained frontier. 
Lower IGD values indicate that the solutions are closer to the reference 
Pareto front, reflecting better convergence.

\textbf{Spread} measures how evenly distributed the solutions are along the 
Pareto frontier. In this study, spread is computed as the mean pairwise 
distance between solutions, providing an indication of diversity across 
the objective space.

\textbf{Mean Distance to Heaven (MDH)} measures the distance of solutions 
to an ideal point, where all objectives are optimized. Lower MDH values 
indicate solutions that are closer to the ideal trade-off.

All objectives are normalized prior to evaluation. Hypervolume is computed 
using a reference point slightly outside the normalized space, which may 
result in values greater than 1. These values remain comparable within 
each dataset.

\subsection{A Gap in the Literature}

To better understand how prior work evaluates multi-objective optimization 
algorithms, we conducted a literature search in March 2026 using the query 
``multi-objective optimization'' in Google Scholar. This broad query was 
chosen to capture influential papers across different domains without 
restricting results to specific algorithms or applications.

The top 100 papers were collected and ranked by citation count. Following 
common bibliometric practice, we identified the ``knee'' of the citation 
curve, defined as the point furthest from the line connecting the most and 
least cited papers. The knee occurred at rank 13, corresponding to papers 
with at least 627 citations. These highly cited papers were then analyzed 
in detail to identify common patterns in the literature.

\subsection{Citation Analysis}

\begin{figure}[h]
\centering
\includegraphics[width=0.8\linewidth]{knee_graph.png}
\caption{Citation counts of the top 100 papers ranked by citation count. The knee occurs at rank 13.}
\end{figure}

The knee point represents a natural cutoff between the most influential 
papers and the remainder of the literature. We focus on these above-knee 
papers to understand dominant research directions.

\subsection{Classification of Prior Work}

From the above-knee papers, we identified five key properties used to 
characterize prior work:

\begin{itemize}
    \item \textbf{Optimization framework:} Some studies use frameworks such 
    as pymoo~\cite{blank2020pymoo} and PlatEMO~\cite{tian2017platemo}, which 
    allow multiple algorithms to be evaluated within a single environment. 
    However, these studies focus on comparing individual algorithms rather 
    than combining their outputs.

    \item \textbf{SE datasets:} Most studies rely on synthetic benchmark 
    functions or domain-specific datasets outside software engineering, 
    such as energy optimization~\cite{cui2017energy} or wireless sensor 
    networks~\cite{fei2016wsn}. Very few evaluate algorithms on diverse 
    software engineering datasets.

    \item \textbf{$\geq$5 optimization problems:} Many studies evaluate 
    algorithms on only a small number of problems. While some benchmarking 
    frameworks consider larger suites, these are typically synthetic rather 
    than real-world SE tasks.

    \item \textbf{Performance metrics:} Metrics such as hypervolume and 
    spread are widely used~\cite{riquelme2015metrics,emmerich2018tutorial}, 
    but they are rarely applied together across a large and diverse set of 
    problems.

    \item \textbf{Algorithm ensembles:} Very few studies combine multiple 
    multi-objective algorithms. Some work applies MOEAs to ensemble learning 
    in machine learning~\cite{chandra2006ensemble,minku2013effort}, and one 
    study combines multiple MOEAs on synthetic benchmarks~\cite{zhou2017ensemble}, 
    but none evaluates ensembles on software engineering datasets.
\end{itemize}

\subsection{Venn Analysis of Prior Work}

\begin{figure}[h]
\centering
\includegraphics[width=0.6\linewidth]{venn_diagram.png}
\caption{Overlap of five key properties in prior work. No paper satisfies all properties simultaneously.}
\end{figure}

The overlap of these five properties reveals a clear gap in the literature. 
No prior work satisfies all five properties simultaneously, particularly 
the combination of algorithm ensembles evaluated on diverse software 
engineering datasets.

\subsection{Comparison of Prior Work}

\begin{table*}[h]
\centering
\caption{Summary of above-knee papers and their properties.}
\begin{tabular}{lcccccc}
\hline
Paper & Framework? & SE Data? & $\geq$5 Problems? & Metrics? & Ensemble? \\
\hline
Deb (2016) & $\times$ & $\times$ & -- & -- & $\times$ \\
Blank (2020) & $\checkmark$ & $\times$ & $\checkmark$ & $\checkmark$ & $\times$ \\
Tian (2017) & $\checkmark$ & $\times$ & $\checkmark$ & $\checkmark$ & $\times$ \\
Sener (2018) & $\times$ & $\times$ & -- & -- & $\times$ \\
Gunantara (2018) & $\times$ & $\times$ & $\times$ & $\checkmark$ & $\times$ \\
Emmerich (2018) & $\times$ & $\times$ & $\times$ & $\checkmark$ & $\times$ \\
Mirjalili (2018) & $\times$ & $\times$ & $\checkmark$ & $\checkmark$ & $\times$ \\
Mirjalili (2017) & $\times$ & $\times$ & $\checkmark$ & $\checkmark$ & $\times$ \\
Cui (2017) & $\times$ & $\times$ & -- & -- & $\times$ \\
Li (2018) & $\times$ & $\times$ & $\checkmark$ & $\checkmark$ & $\times$ \\
Riquelme (2015) & $\times$ & $\times$ & -- & $\checkmark$ & $\times$ \\
Tian (2020) & $\times$ & $\times$ & -- & $\checkmark$ & $\times$ \\
Fei (2016) & $\times$ & $\times$ & -- & $\checkmark$ & $\times$ \\
\textbf{This Work} & $\checkmark$ & $\checkmark$ & $\checkmark$ & $\checkmark$ & $\checkmark$ \\
\hline
\end{tabular}
\end{table*}

From this analysis, we observe that no prior work satisfies all five 
properties simultaneously. In particular, no study combines multiple 
multi-objective algorithms as an ensemble while evaluating them on a 
diverse set of software engineering datasets. This gap motivates the 
proposed study.

\section{Motivation}

Although many optimization algorithms have been proposed, selecting the 
best algorithm for a particular problem remains challenging. Each 
algorithm uses different strategies to explore the search space and 
maintain diversity among candidate solutions.

Because of these differences, algorithms may discover different parts 
of the Pareto frontier. For example, one algorithm may perform well in 
regions that emphasize one objective, while another algorithm may 
identify solutions that prioritize different objectives.

If these algorithms produce complementary results, combining their 
solutions could potentially lead to a more complete approximation of 
the Pareto frontier. This idea is similar to ensemble learning in machine learning, where multiple 
models are combined to improve overall prediction performance~\cite{dietterich2000ensemble}. 
Similar concepts have also been explored in multi-objective optimization~\cite{chandra2006ensemble}.

Despite the success of ensemble methods in other areas of computer 
science, relatively little research has explored ensemble techniques 
in multi-objective optimization, particularly for software engineering 
datasets. This gap motivates the investigation conducted in this study.

\section{Methods}

\subsection{Variables}

The independent variable in this study is the optimization algorithm 
used to generate Pareto frontiers. Multiple algorithms available in 
the pymoo framework will be evaluated.

The dependent variables are performance metrics used to measure the 
quality of the resulting Pareto frontiers. These metrics include 
hypervolume and spread, which measure convergence and diversity.

\subsection{Dataset Selection}

Datasets used in this study are obtained from the MOOT repository, 
which contains a collection of multi-objective optimization tasks 
derived from software engineering problems. These datasets include 
problems related to system configuration, hyperparameter tuning, 
and project decision-making tasks.

The use of real software engineering datasets ensures that the 
experimental results reflect practical optimization challenges 
rather than purely synthetic benchmark problems.

For this study, we used 69 datasets after applying consistency filters: 
each dataset must contain at least 1000 rows and at least 2 objective 
columns (identified by the MOOT +/- objective naming convention).


\subsection{Experimental Methodology}

The experimental process is structured into three stages, referred to 
as Pass 1, Pass 2, and Pass 3. Each stage builds upon the previous one 
to evaluate the effectiveness of ensemble-based optimization.

\subsubsection{Pass 1: Individual Optimization}

In the first stage, each optimization algorithm is run independently 
on each dataset. The output of each run is a Pareto frontier representing 
a set of non-dominated solutions.

We evaluate nine multi-objective algorithms from \textit{pymoo}: NSGA-II, 
MOEA/D, NSGA-III, UNSGA-III, RNSGA-III, RVEA, AGEMOEA, AGEMOEA2, and CTAEA.

This stage allows us to evaluate how different algorithms perform 
individually and to observe differences in the regions of the solution 
space that each algorithm explores.

\subsubsection{Pass 2: Ensemble Frontier Construction}

In the second stage, the Pareto frontiers produced by all algorithms 
are combined into a single set of candidate solutions. From this combined 
set, dominated solutions are removed to produce an ensemble Pareto frontier.

This ensemble frontier represents the union of solutions discovered 
by multiple algorithms and is intended to provide a more complete 
approximation of the Pareto front.

\subsubsection{Pass 3: Ensemble-Derived Initialization}

In the third stage, the Pareto frontiers produced by multiple algorithms 
are combined and dominated solutions are removed to form a unified 
ensemble frontier. These solutions are then used as seed data to 
initialize a single optimization algorithm.

Unlike Pass 1, which begins with random initialization, this stage 
leverages solutions discovered in earlier stages. This allows us to 
evaluate whether ensemble-derived initialization can improve 
optimization performance compared to standard random initialization.

\section{Data Analysis}

The performance of each algorithm and ensemble configuration will be 
evaluated using commonly used multi-objective optimization metrics.

Hypervolume will be used to measure how much of the objective space is 
covered by the Pareto frontier. Higher hypervolume values indicate that 
the solutions dominate a larger region of the objective space and are 
closer to the optimal frontier.

Spread will be used to evaluate the diversity of solutions along the 
Pareto frontier. A good Pareto frontier should include solutions that 
cover a wide range of trade-offs across the objective space.

Inverted Generational Distance (IGD) will be used to measure how close 
the obtained frontier is to a reference frontier. Lower IGD indicates 
better convergence toward the best known trade-off surface.

Mean Distance to Heaven (MDH) will be used to measure how close solutions 
are to the ideal point in objective space. Lower MDH indicates better 
overall objective quality relative to the ideal target.

By comparing these metrics across multiple algorithms and ensemble 
configurations, we can determine whether combining Pareto frontiers 
improves optimization performance on software engineering datasets.


\section{Results}

The experimental results show that no single optimization algorithm 
consistently outperforms all others across all datasets and evaluation 
metrics. Instead, different algorithms perform best depending on the 
specific metric being considered.

The percentages reported in this section are computed as pairwise win rates 
against Pass~1 results. For a metric $m$, we define:
\[
\text{Better\%}_m = 100 \times \frac{W_m}{N_m},
\]
where $W_m$ is the number of valid pairwise comparisons in which the evaluated 
method is better than the corresponding Pass~1 result, and $N_m$ is the total 
number of valid pairwise comparisons for metric $m$. For HV and Spread, 
``better'' means higher value; for IGD and MDH, ``better'' means lower value.

Table~\ref{tab:pass2-vs-pass1} summarizes Pass~2 (ensemble frontier) against 
Pass~1 (individual algorithms), using percentage of pairwise wins over all 
dataset-algorithm comparisons. Pass~2 is strongly better on HV ($80.27\%$) 
and IGD ($91.57\%$), and also better on Spread ($62.08\%$) and MDH ($50.77\%$). 
This indicates that combining fronts improves both convergence-oriented metrics 
and diversity for a majority of comparisons.

\begin{table}[h]
\centering
\caption{Pass~2 vs Pass~1: Percent of comparisons where Pass~2 is better (HV/Spread higher; IGD/MDH lower).}
\label{tab:pass2-vs-pass1}
\begin{tabular}{lcccc}
\hline
Method & HV (\%) & Spread (\%) & IGD (\%) & MDH (\%) \\
\hline
Pass~2 Ensemble & 80.27 & 62.08 & 91.57 & 50.77 \\
\hline
\end{tabular}
\end{table}

For Pass~3 (ensemble-derived initialization), Table~\ref{tab:pass3-vs-pass1} 
reports the same comparison per algorithm. No single algorithm dominates all 
metrics: RVEA has the highest HV ($87.18\%$) and MDH ($73.68\%$) win rates, 
AGEMOEA has the highest IGD win rate ($84.62\%$), and RNSGA3 has the highest 
Spread win rate ($83.45\%$). This reinforces that different algorithms provide 
different strengths after ensemble-based refinement.

\begin{table*}[h]
\centering
\caption{Pass~3 vs Pass~1: Percent of comparisons where each Pass~3 algorithm is better (HV/Spread higher; IGD/MDH lower).}
\label{tab:pass3-vs-pass1}
\begin{tabular}{lcccc}
\hline
Algorithm & HV (\%) & Spread (\%) & IGD (\%) & MDH (\%) \\
\hline
AGEMOEA  & 74.02 & 67.48 & \textbf{84.62} & 40.51 \\
AGEMOEA2 & 72.48 & 65.61 & 84.44 & 41.54 \\
CTAEA    & 73.23 & 69.06 & 84.01 & 47.81 \\
MOEAD    & 44.05 & 42.27 & 42.26 & 53.97 \\
NSGA2    & 70.37 & 70.53 & 83.84 & 43.43 \\
NSGA3    & 69.57 & 73.27 & 78.80 & 49.06 \\
RNSGA3   & 32.69 & \textbf{83.45} & 35.68 & 41.03 \\
RVEA     & \textbf{87.18} & 40.30 & 81.20 & \textbf{73.68} \\
UNSGA3   & 69.91 & 73.83 & 79.66 & 50.77 \\
\hline
\end{tabular}
\end{table*}

Overall, these results support the main hypothesis that ensemble information 
is useful: Pass~2 provides strong gains across all four metrics, and Pass~3 
further improves outcomes in algorithm-specific ways rather than with a single 
universally best optimizer.

\section{Discussion}

The results demonstrate that multi-objective optimization algorithms 
exhibit different strengths depending on the evaluation metric and 
dataset. No single algorithm consistently dominates across all 
experiments, which aligns with the No Free Lunch theorem for optimization~\cite{wolpert1997nfl}.

This variability suggests that algorithms are not redundant, but rather 
complementary. Each algorithm tends to discover different regions of the 
Pareto frontier, emphasizing different trade-offs among objectives.

The ensemble method benefits from this complementarity by combining 
solutions from multiple algorithms. By pooling Pareto frontiers and 
removing dominated solutions, the ensemble is able to construct a more 
complete approximation of the solution space.

However, the effectiveness of the ensemble approach varies across 
datasets. While some datasets show clear improvements, others show only 
minor gains. This indicates that the success of ensemble optimization 
may depend on factors such as dataset size, dimensionality, and the 
complexity of the objective landscape.

The refinement stage further highlights this behavior. While 
initialization using ensemble-derived solutions can improve performance 
in some cases, it does not universally outperform standard initialization 
methods.

\section{Threats to Validity}

Several limitations should be considered when interpreting the results 
of this study.

First, the reference Pareto front used for evaluation is approximated 
from available data and may not represent the true optimal frontier. 
This may affect metrics such as IGD.

Second, the spread metric used in this study is implemented as the mean 
pairwise distance between solutions, which differs from standard spread 
metrics used in the literature.

Third, the refinement stage operates only on solutions discovered during 
earlier optimization stages and does not explore the full search space. 
This limits its ability to discover entirely new regions of the Pareto 
frontier.

Finally, the results may be sensitive to parameter choices such as 
population size, number of generations, and random seed. Different 
parameter settings could lead to different outcomes.

\section{Future Work}

A natural next step is to optimize each individual algorithm used in 
Pass~1 with automated hyperparameter optimization methods such as SMAC. 
In this study, all algorithms were run with a common default configuration, 
which supports fair comparison but may not reflect each algorithm's best 
possible performance on each dataset. Per-algorithm tuning could produce 
stronger Pass~1 frontiers, which would in turn provide higher-quality 
inputs to the ensemble stages.

Improving the quality of Pass~1 frontiers is likely to improve Pass~3 as 
well, since Pass~3 depends on the combined frontier produced by earlier 
stages. Better initial fronts can improve both convergence and diversity 
during refinement, potentially yielding stronger final Pareto solutions.

Another direction is to separate algorithm selection decisions across stages. 
The best set of algorithms for Pass~1 (maximizing complementary frontier 
coverage) may not be the same as the best choice for Pass~3 refinement 
(maximizing improvement given seeded initialization). Future work can study 
stage-specific portfolios, where one algorithm subset is chosen for Pass~1 
and a different subset is chosen for Pass~3 based on metric-specific goals.

\section{Conclusion}

This paper investigated the use of ensemble techniques for multi-objective 
optimization on software engineering datasets. The results show that no 
single algorithm consistently dominates across all metrics and datasets. 
Instead, different algorithms excel in different aspects of optimization, 
indicating that they explore distinct regions of the solution space.

By combining Pareto frontiers from multiple algorithms, the ensemble 
approach leverages this complementarity to produce improved results in 
several cases. While the ensemble does not universally outperform all 
individual algorithms, it often provides a more balanced and robust 
approximation of the Pareto frontier.

These findings suggest that ensemble-based optimization is a promising 
direction for future research in search-based software engineering. 
Future work could explore adaptive ensemble methods and investigate how 
dataset characteristics influence the effectiveness of ensemble 
approaches.


\bibliographystyle{ACM-Reference-Format}
\bibliography{references}

\end{document}